\section{Introduction}

Static Application Security Testing (SAST) tools are a cornerstone of secure
software development, yet their practical value is undermined by noise.
The Ghost Security CAST report found a 99.5\% false positive rate for
command-injection findings in Python/Flask applications, with 91\% of all
2{,}166 findings across categories classified as noise~\cite{ghost2025cast}.
The picture is similar at scale: NIST's SATE~V and SATE~VI evaluations
measured false positive rates of 70--92\% depending on language and
tool~\cite{delaitre2018satev,delaitre2023satevi}. On the detection side,
the Fluid Attacks benchmark showed that the best fully automated tool found
only 22.7\% of known vulnerabilities, with an average F-score of 1.9\%
across 36 commercial and open-source scanners~\cite{fluidattacks2025}.
Practitioners, faced with hundreds of alerts they must manually triage,
increasingly ignore scanner output altogether.

A new generation of LLM-powered security scanners promises to change this.
By leveraging large language models for semantic code understanding rather than
pattern matching, these tools claim deeper reasoning about data flow, context,
and exploitability. Yet no open benchmark exists to verify whether LLM-based
scanners actually deliver on these promises---or whether they simply trade one
set of failure modes for another. Existing benchmarks fall short on one or more
critical dimensions:

\begin{itemize}
  \item \textbf{Synthetic code.} The OWASP Benchmark~\cite{owaspbenchmark}
    consists of synthetic Java test cases that do not reflect the structure,
    complexity, or idioms of real-world applications.
  \item \textbf{Closed data.} SastBench~\cite{feiglin2026sastbench} is the
    closest prior work to ours, but its dataset has not been publicly released
    and its evaluation focuses on triage agents rather than scanner accuracy.
  \item \textbf{Vendor self-assessment.} Benchmarks published by
    ZeroPath~\cite{zeropath2024}, Cycode~\cite{cycode2024benchmark}, and
    DryRun~\cite{dryrun2025sast} evaluate their own tools on proprietary
    datasets, making independent reproduction impossible.
  \item \textbf{No false positive measurement.} Most benchmarks report
    detection rates alone, ignoring the false positive burden that dominates
    practitioner experience.
\end{itemize}

\noindent
We expand on these limitations in Section~\ref{sec:related-work}. The gap is
clear: the community lacks a reproducible, open benchmark that tests both
traditional and LLM-based scanners on real code, measuring not just what they
find but also what they falsely flag.

This paper makes three contributions:

\begin{enumerate}
  \item \textbf{RealVuln}, the first fully open-source SAST benchmark built
    from real-world code. It comprises 26 intentionally vulnerable Python
    repositories with 796 hand-labeled findings---including 120 false positive
    traps---scored using an F3-weighted metric ($\beta{=}3$, recall weighted 9$\times$ over precision) as the primary metric for production-grade security assessment, with F2 ($\beta{=}2$) reported as a secondary metric for lower-risk contexts.
  \item \textbf{The first large-scale head-to-head comparison} of 15 security
    scanners on identical ground truth: 3 traditional SAST tools (Semgrep,
    Snyk, SonarQube), 11 LLM-based scanner configurations spanning 7 model families, and
    1 proprietary hybrid system.
  \item \textbf{A living benchmark} with versioned scoring, a public
    dashboard, and a clear contribution path so that the community can add new
    scanners, repositories, and ground-truth labels over time.
\end{enumerate}

We release all code, ground-truth data, scanner results, scoring tooling, and
the interactive dashboard at
\url{https://github.com/kolega-ai/RealVulnBenchmark}.
